\section{Избежание стимулов к вмешательству в глубоких алгоритмах обучения с подкреплением с помощью метода раздельного одобрения}

\subsection{Ключевые идеи}

Цель статьи~\cite{uesato2020avoidingtamperingincentivesdeep}~-- разработать методы обучения с подкреплением, при которых агент не имеет стимула вмешиваться в источник обратной связи.  
Для описания такой ситуации вводится модель марковского процесса принятия решений с искажённой обратной связью, где наблюдаемое одобрение $\tilde{D}$ может отличаться от истинного $D$ из-за функции искажения $c(s, k, d)$.

Предлагается метод раздельного одобрения, в котором агент получает оценку не за выполненное действие, а за другое, независимое действие, выбранное тем же распределением вероятностей.  
Разделение действия и запроса на одобрение делает обратную связь независимой от поведения агента в среде и исключает стимулы к её искажению. То есть агент оптимизирует стратегию по реальному качеству действий, а не по результатам, которые можно искусственно завысить, воздействуя на канал обратной связи.

\subsection{Математическое описание метода}

Среда формализуется как марковский процесс принятия решений с искажённой обратной связью:
\begin{equation}
\mu = (S, A, D, K, p, f, r, \gamma, c),
\end{equation}
где $S$~-- множество состояний, $A$~-- действий, $D$~-- оценок, $K$~-- запросов, $p$~-- начальное распределение, $f$~-- функция переходов, $r$~-- истинное вознаграждение, $\gamma$~-- коэффициент дисконтирования, $c$~-- функция искажения обратной связи.

Искажение описывается как функция $c(s, k, d)$, которая изменяет истинное одобрение $d$. В частном случае аддитивного, независимого от запроса искажения:
\begin{equation}
c(s, k, d) = d + c_s.
\end{equation}
Тогда наблюдаемое агентом одобрение:
\begin{equation}
\tilde{D}_{t+1} = c(S_{t+1}, K_t, D_{t+1}).
\end{equation}

\paragraph{Метод с градиентом политики.} Агент обновляет параметры стратегии по наблюдаемому одобрению:
\begin{equation}
\theta_{t+1} = \theta_t + \alpha\, \tilde{d}\, \nabla_\theta \log \pi_\theta(k|s),
\end{equation}
где $\pi_\theta(a|s)$~-- вероятность выбора действия $a$ в состоянии $s$.
Среднее направление обновления совпадает с направлением в "честной" среде:
\begin{equation}
\mathbb{E}_{\mu,\pi_\theta}[c(s',k,\delta(s,k))\nabla_\theta \log \pi_\theta(k|s)] =
\mathbb{E}_{\bar{\mu},\pi_\theta}[\delta(s,a)\nabla_\theta \log \pi_\theta(a|s)].
\end{equation}
Это означает, что обучение в среде, подверженной вмешательству, эквивалентно обучению в идеальной среде без искажений.

\paragraph{Метод на основе функции качества.} Для функции качества $Q(s,k)$ обновление выполняется:
\begin{equation}
Q_{t+1}(s,k) = (1 - \alpha) Q_t(s,k) + \alpha\, \tilde{d},
\end{equation}
где шаг обучения корректируется:
\begin{equation}
\alpha = \alpha_{начальное}(M_t(s)\pi_K(k|s))^{-1}.
\end{equation}

\subsection{Алгоритм}

Первый метод DA-PG (Градиенты политики с разъединенным одобрением) расширяет стандартные градиенты политики путем независимого сэмплирования действия для выполнения в среде и действия для запроса обратной связи, что устраняет корреляцию между вмешательством и получаемой обратной связью.

\begin{enumerate}
    \item Инициализировать параметры политики $\theta_0$
    \item Для каждого шага $t = 0$ до $T$:
    \begin{itemize}
        \item Наблюдать текущее состояние $s$
        \item Независимо сэмплировать действие $a \sim \pi_{\theta_t}(s)$ и запрос $k \sim \pi_{\theta_t}(s)$
        \item Получить следующее состояние $s' \sim f(s,a)$ и искаженное одобрение $\tilde{d} = c(s',k,\delta(s,k))$
        \item Обновить параметры политики: $\theta_{t+1} := \theta_t + \alpha\tilde{d}\nabla_\theta \log \pi_\theta(k \mid s)$
    \end{itemize}
\end{enumerate}

Второй метод DA-QL (Q-обучение с разъединенным одобрением) расширяет стандартное Q-обучение аналогичным принципом разъединения, но с добавлением поправки на важность сэмплирования для компенсации смещения в обновлениях Q-значений.

\begin{enumerate}
    \item Инициализировать Q-значения $Q_0(s,k) = 0$ и счетчики посещений $M_0(s)$
    \item Для каждого шага $t = 0$ до $T$:
    \begin{itemize}
        \item Наблюдать текущее состояние $s$
        \item Определить $\epsilon$-жадные политики $\pi_A$ и $\pi_K$ с разными параметрами исследования
        \item Независимо сэмплировать действие $a \sim \pi_A(s)$ и запрос $k \sim \pi_K(s)$
        \item Вычислить learning rate: $\alpha := \alpha_{\text{init}}(M_t(s)\pi_K(k \mid s))^{-1}$
        \item Получить следующее состояние $s' \sim f(s,a)$ и искаженное одобрение $\tilde{d} = c(s',k,\delta(s,k))$
        \item Обновить Q-значение: $Q_{t+1}(s,k) := (1-\alpha)Q_t(s,k) + \alpha\tilde{d}$
        \item Обновить счетчик посещений: $M_{t+1}(s') := M_t(s') + 1$
    \end{itemize}
\end{enumerate}


\subsection{Наборы данных}

Эксперименты проводились в среде REALab~\cite{Kumar2020Realab}~-- трёхмерном симуляторе, предназначенном для исследования устойчивости искажённой обратной связи. Среда позволяет агенту не только взаимодействовать с объектами, но и изменять элементы, определяющие вычисление вознаграждения, что моделирует возможность вмешательства в источник обратной связи.

Основная задача~-- \emph{разблокировать дверь}, где агент должен собирать яблоки и открывать дверь для получения вознаграждения. При этом он имеет доступ к физическим компонентам, влияющим на механизм оценки, что создаёт условия для изучения поведения при наличии стимулов к вмешательству.

Дополнительно рассматривались упрощённые сценарии: Поиск и избегание (Seek-Avoid), а также табличные модели с явно заданной функцией искажения обратной связи. Эти сценарии использовались для проверки обобщаемости и устойчивости метода раздельного одобрения.

\subsection{Метрики}

Для количественной оценки эффективности и устойчивости алгоритмов использовались следующие метрики:

\begin{itemize}
  \item \emph{Общий возврат} (Return)~-- ожидаемое суммарное вознаграждение, полученное агентом за эпизод согласно исходной функции вознаграждения среды. Эта метрика отражает реальную производительность агента в решении поставленной задачи.

  \item \emph{Процент эпизодов с вмешательством} (Tampering \% Episodes)~-- доля эпизодов, в которых разница между истинной и наблюдаемой обратной связью превышает порог $10.0$ хотя бы в один момент времени. Эпизод классифицируется как содержащий вмешательство, если $|\tilde{D}_t - D_t| > 10.0$ для любого $t$ в эпизоде.
\end{itemize}


\subsection{Модели}

Для реализации и сравнения методов использовались следующие архитектуры и подходы:

\begin{itemize}
  \item \emph{R2D2 (Recurrent Replay Distributed DQN)}~\cite{kapturowski2019recurrent}~-- архитектура с рекуррентной нейронной сетью и распределенным воспроизведением опыта, использовалась как основа для Q-обучения.
  
  \item \emph{PPO (Proximal Policy Optimization)}~\cite{schulman2017proximal}~-- метод проксимальной оптимизации политики, применявшийся для алгоритмов градиента политики с гиперпараметрами по умолчанию из оригинальной статьи.
  
  \item \emph{REALab}~\cite{kumar2020realab}~-- платформа для симулированных сред со встроенной обратной связью, используемая для моделирования проблемы вмешательства.
  
  \item \emph{Сравнительные подходы}:
  \begin{itemize}
    \item Standard RL (R2D2 с оптимизацией наблюдаемых вознаграждений)
    \item Approval QL/PG (одобрение без разъединения)
    \item Tampering Detector (стандартный RL с детектором вмешательства)
    \item Offline Imitation и DAgger~\cite{ross2011reduction} (имитационное обучение)
    \item Non-embedded (агент с оракульным доступом к истинным вознаграждениям)
  \end{itemize}
\end{itemize}

\subsection{Результаты}

Экспериментальные результаты, представленные в таблице~\ref{fig:table},  Cтандартные алгоритмы RL демонстрируют высокий уровень вмешательства (78\% для Standard RL), в то время как предложенные методы разъединенного одобрения значительно снижают этот показатель (6\% для DA-PG и 12\% для DA-QL).

\begin{table}[h]
\centering
\begin{tabular}{lcc}
\hline
\textbf{Алгоритм} & \textbf{Вмещательство (\% эпизодов)} & \textbf{Возврат (средний)} \\
\hline
Standard RL & 78\% & 1.0 \\
Tampering Detector & 84\% & 1.2 \\
Non-embedded & 6\% & 14.8 \\
\hline
Approval QL & 26\% & 9.5 \\
DA-QL (без поправки) & 19\% & 10.2 \\
DA-QL & 12\% & 11.9 \\
\hline
Approval PG & 18\% & 13.4 \\
DA-PG & 6\% & 13.7 \\
\hline
Offline Imitation & 8\% & 14.0 \\
DAgger & 8\% & 14.0 \\
\hline
\end{tabular}
\caption{Сравнение методов по уровню вмешательства и производительности в задаче "Открыть Дверь"}
\label{fig:table}
\end{table}
